\documentclass[11pt,a4paper]{article}

\usepackage[utf8]{inputenc}
\usepackage[T1]{fontenc}
\usepackage[margin=2.5cm]{geometry}
\usepackage{hyperref}
\usepackage{listings}
\usepackage{xcolor}
\usepackage{booktabs}
\usepackage{caption}
\usepackage{microtype}
\usepackage{parskip}
\usepackage{titlesec}
\usepackage{enumitem}
\usepackage{lmodern}

% --- Python listings style ---
\definecolor{codebg}{RGB}{248,248,248}
\definecolor{keyword}{RGB}{0,0,180}
\definecolor{string}{RGB}{163,21,21}
\definecolor{comment}{RGB}{0,128,0}
\definecolor{annotation}{RGB}{130,65,0}

\lstdefinelanguage{Python}{
  keywords={def,class,return,if,elif,else,match,case,from,import,pass,None,True,False,and,or,not,in,is,for,while,with,as,try,except,finally,raise,lambda,yield,del,assert,break,continue,global,nonlocal},
  keywordstyle=\color{keyword}\bfseries,
  sensitive=true,
  comment=[l]{\#},
  commentstyle=\color{comment}\itshape,
  stringstyle=\color{string},
  morestring=[b]",
  morestring=[b]',
}

\lstset{
  language=Python,
  basicstyle=\ttfamily\small,
  backgroundcolor=\color{codebg},
  frame=single,
  rulecolor=\color{gray!40},
  framesep=4pt,
  xleftmargin=6pt,
  xrightmargin=6pt,
  showstringspaces=false,
  tabsize=4,
  breaklines=true,
  captionpos=b,
}

% --- Section formatting ---
\titleformat{\section}{\large\bfseries}{}{0em}{\thesection\quad}
\titleformat{\subsection}{\normalsize\bfseries}{}{0em}{\thesubsection\quad}

\hypersetup{
  colorlinks=true,
  linkcolor=blue!60!black,
  urlcolor=blue!60!black,
}

% --- Inline code ---
\newcommand{\code}[1]{\texttt{#1}}

\begin{document}

% -----------------------------------------------------------------------
\begin{titlepage}
  \centering
  \vspace*{3cm}
  {\LARGE\bfseries Technical Report\\[0.5em]
   Python Refactoring Detection in RefactoringMiner\par}
  \vspace{1.5cm}
  {\large Type Hint Annotations and Structural Pattern Matching\par}
  \vspace{2cm}
  {\normalsize March 2026\par}
  \vfill
\end{titlepage}

% -----------------------------------------------------------------------
\tableofcontents
\newpage

% -----------------------------------------------------------------------
\section{Overview}

This report describes two sets of new refactoring detection capabilities added to RefactoringMiner for Python source code. The changes were implemented across two sessions and introduce eleven new refactoring types: nine related to \textbf{PEP 484/526 type hint annotations} and two related to \textbf{structural pattern matching} (Python~3.10+).

% -----------------------------------------------------------------------
\section{Motivation}

\subsection{Type Hint Annotations}

Python introduced optional type hints via \href{https://peps.python.org/pep-0484/}{PEP~484} (Python~3.5) and variable annotations via \href{https://peps.python.org/pep-0526/}{PEP~526} (Python~3.6). Since then, two major forces have driven annotation-related changes in real-world Python codebases.

\begin{enumerate}[leftmargin=*, label=\arabic*.]
  \item \textbf{Gradual typing adoption.} Teams incrementally annotate previously unannotated functions and variables as they mature their codebase or integrate static analysis tools such as mypy, pyright, or Pytype.

  \item \textbf{Modernisation from the \code{typing} module to built-in generics.}
        \href{https://peps.python.org/pep-0585/}{PEP~585} (Python~3.9) allowed \code{list[int]} instead of \code{typing.List[int]}, and \href{https://peps.python.org/pep-0604/}{PEP~604} (Python~3.10) introduced the \code{X | Y} union syntax as a replacement for \code{typing.Optional[X]} and \code{typing.Union[X, Y]}. Many projects have been migrating to these shorter, import-free forms.
\end{enumerate}

RefactoringMiner already detected analogous changes in Java (e.g., \emph{Change Parameter Type}). Extending detection to Python type annotations closes this gap and enables empirical studies of typing evolution in Python projects.

\subsection{Structural Pattern Matching}

Python~3.10 introduced structural pattern matching (\code{match}/\code{case}) via \href{https://peps.python.org/pep-0634/}{PEP~634}. It provides a more expressive alternative to long \code{if}/\code{elif}/\code{else} chains when dispatching on the shape or value of data. Developers may migrate in either direction:

\begin{itemize}[leftmargin=*]
  \item \textbf{Forward:} replace a chain of equality comparisons with a \code{match} statement when the logic maps cleanly to case patterns.
  \item \textbf{Backward:} revert a \code{match} statement to a conditional chain, for instance when targeting older Python interpreters or when the conditional form is judged more readable.
\end{itemize}

Detecting both directions allows researchers and tool authors to study adoption patterns of this language feature in the wild.

% -----------------------------------------------------------------------
\section{Implemented Refactoring Types}

\subsection{Type Hint Annotation Refactorings}

Table~\ref{tab:annotation-types} lists the nine new \code{RefactoringType} constants introduced for type hint annotations.

\begin{table}[ht]
  \centering
  \caption{New refactoring types for Python type hint annotations.}
  \label{tab:annotation-types}
  \small
  \begin{tabular}{@{}ll@{}}
    \toprule
    \textbf{Refactoring Type} & \textbf{Trigger} \\
    \midrule
    \code{ADD\_PARAMETER\_TYPE\_ANNOTATION}    & A parameter that had no type hint gains one \\
    \code{REMOVE\_PARAMETER\_TYPE\_ANNOTATION} & A parameter loses its type hint \\
    \code{CHANGE\_PARAMETER\_TYPE\_ANNOTATION} & A parameter's type hint changes \\
    \code{ADD\_RETURN\_TYPE\_ANNOTATION}        & A function gains a \code{-> T} return annotation \\
    \code{REMOVE\_RETURN\_TYPE\_ANNOTATION}     & A function loses its \code{-> T} return annotation \\
    \code{CHANGE\_RETURN\_TYPE\_ANNOTATION}     & A function's return annotation changes \\
    \code{ADD\_VARIABLE\_TYPE\_ANNOTATION}      & A local variable or class attribute gains \code{: T} \\
    \code{REMOVE\_VARIABLE\_TYPE\_ANNOTATION}   & A local variable or class attribute loses its annotation \\
    \code{CHANGE\_VARIABLE\_TYPE\_ANNOTATION}   & A variable's annotation changes \\
    \bottomrule
  \end{tabular}
\end{table}

\subsection{Structural Pattern Matching Refactorings}

Table~\ref{tab:pattern-types} lists the two new \code{RefactoringType} constants for structural pattern matching.

\begin{table}[ht]
  \centering
  \caption{New refactoring types for structural pattern matching.}
  \label{tab:pattern-types}
  \small
  \begin{tabular}{@{}ll@{}}
    \toprule
    \textbf{Refactoring Type} & \textbf{Trigger} \\
    \midrule
    \code{REPLACE\_CONDITIONAL\_WITH\_PATTERN\_MATCHING} & An \code{if}/\code{elif}/\code{else} chain is replaced by \code{match}/\code{case} \\
    \code{REPLACE\_PATTERN\_MATCHING\_WITH\_CONDITIONAL} & A \code{match}/\code{case} statement is replaced by \code{if}/\code{elif}/\code{else} \\
    \bottomrule
  \end{tabular}
\end{table}

% -----------------------------------------------------------------------
\section{Key Examples}

\subsection{Add Parameter and Return Type Annotations}

A function with no type information gains full annotations --- a typical first step when adopting static analysis tooling.

\begin{lstlisting}[caption={Before: unannotated function.}]
class Calculator:
    def add(self, x, y):
        return x + y
\end{lstlisting}

\begin{lstlisting}[caption={After: fully annotated function.}]
class Calculator:
    def add(self, x: int, y: int) -> int:
        return x + y
\end{lstlisting}

Detected refactorings:
\begin{itemize}[leftmargin=*]
  \item \code{ADD\_PARAMETER\_TYPE\_ANNOTATION} for \code{x} (added \code{: int})
  \item \code{ADD\_PARAMETER\_TYPE\_ANNOTATION} for \code{y} (added \code{: int})
  \item \code{ADD\_RETURN\_TYPE\_ANNOTATION} (added \code{-> int})
\end{itemize}

\subsection{Add Return Type Annotation Only}

A function already has parameter annotations but is missing a return annotation --- a common partial-typing scenario.

\begin{lstlisting}[caption={Before: missing return annotation.}]
class DataService:
    def get_name(self):
        return ""
\end{lstlisting}

\begin{lstlisting}[caption={After: return annotation added.}]
class DataService:
    def get_name(self) -> str:
        return ""
\end{lstlisting}

Detected refactoring:
\begin{itemize}[leftmargin=*]
  \item \code{ADD\_RETURN\_TYPE\_ANNOTATION} (added \code{-> str})
\end{itemize}

\subsection{Remove All Type Annotations}

The reverse: all annotations are stripped from a function, for example when removing a static analysis dependency or simplifying a prototype.

\begin{lstlisting}[caption={Before: fully annotated.}]
class Calculator:
    def add(self, x: int, y: int) -> int:
        return x + y
\end{lstlisting}

\begin{lstlisting}[caption={After: all annotations removed.}]
class Calculator:
    def add(self, x, y):
        return x + y
\end{lstlisting}

Detected refactorings:
\begin{itemize}[leftmargin=*]
  \item \code{REMOVE\_PARAMETER\_TYPE\_ANNOTATION} for \code{x}
  \item \code{REMOVE\_PARAMETER\_TYPE\_ANNOTATION} for \code{y}
  \item \code{REMOVE\_RETURN\_TYPE\_ANNOTATION}
\end{itemize}

\subsection{Change Parameter Type: \texttt{typing} Module to Built-in Generics (PEP~585)}

A common modernisation pattern when dropping support for Python~$<$~3.9.

\begin{lstlisting}[caption={Before: \texttt{typing.List} generics.}]
from typing import List

class Calculator:
    def add(self, x: List[int], y: List[int]) -> List[int]:
        return x + y
\end{lstlisting}

\begin{lstlisting}[caption={After: built-in \texttt{list} generics (PEP 585).}]
class Calculator:
    def add(self, x: list[int], y: list[int]) -> list[int]:
        return x + y
\end{lstlisting}

Detected refactorings:
\begin{itemize}[leftmargin=*]
  \item \code{CHANGE\_PARAMETER\_TYPE\_ANNOTATION} for \code{x}: \code{List[int]} $\to$ \code{list[int]}
  \item \code{CHANGE\_PARAMETER\_TYPE\_ANNOTATION} for \code{y}: \code{List[int]} $\to$ \code{list[int]}
  \item \code{CHANGE\_RETURN\_TYPE\_ANNOTATION}: \code{List[int]} $\to$ \code{list[int]}
\end{itemize}

\subsection{Change Parameter Type: \texttt{Optional[T]} to \texttt{T | None} (PEP~604)}

\begin{lstlisting}[caption={Before: \texttt{Optional} from \texttt{typing}.}]
from typing import Optional

class Processor:
    def process(self, items: Optional[str]):
        pass
\end{lstlisting}

\begin{lstlisting}[caption={After: union syntax (PEP 604).}]
class Processor:
    def process(self, items: str | None):
        pass
\end{lstlisting}

Detected refactoring:
\begin{itemize}[leftmargin=*]
  \item \code{CHANGE\_PARAMETER\_TYPE\_ANNOTATION} for \code{items}: \code{Optional[str]} $\to$ \code{str | None}
\end{itemize}

\subsection{Change Variable Type Annotation --- Class Attribute}

\begin{lstlisting}[caption={Before: class attribute with \texttt{typing.List}.}]
from typing import List

class Shape:
    items: List[str] = []
\end{lstlisting}

\begin{lstlisting}[caption={After: built-in \texttt{list} (PEP 585).}]
class Shape:
    items: list[str] = []
\end{lstlisting}

Detected refactoring:
\begin{itemize}[leftmargin=*]
  \item \code{CHANGE\_VARIABLE\_TYPE\_ANNOTATION} for \code{items}: \code{List[str]} $\to$ \code{list[str]}
\end{itemize}

\subsection{Change Variable Type Annotation --- Local Variable}

\begin{lstlisting}[caption={Before: local variable annotated as \texttt{int}.}]
class DataProcessor:
    def process(self):
        count: int = 0
        count = count + 1
        return count
\end{lstlisting}

\begin{lstlisting}[caption={After: annotation changed to \texttt{float}.}]
class DataProcessor:
    def process(self):
        count: float = 0
        count = count + 1
        return count
\end{lstlisting}

Detected refactoring:
\begin{itemize}[leftmargin=*]
  \item \code{CHANGE\_VARIABLE\_TYPE\_ANNOTATION} for \code{count}: \code{int} $\to$ \code{float}
\end{itemize}

\subsection{Replace Conditional With Pattern Matching --- Simple Value Dispatch}

The canonical motivating case: a chain of equality tests on a single variable is replaced by a \code{match} statement with literal patterns.

\begin{lstlisting}[caption={Before: \texttt{if}/\texttt{elif}/\texttt{else} chain on a string variable.}]
class ShapeProcessor:
    def describe(self, shape):
        if shape == "circle":
            return "A round shape"
        elif shape == "square":
            return "A shape with four equal sides"
        elif shape == "triangle":
            return "A shape with three sides"
        else:
            return "Unknown shape"
\end{lstlisting}

\begin{lstlisting}[caption={After: \texttt{match}/\texttt{case} with literal patterns.}]
class ShapeProcessor:
    def describe(self, shape):
        match shape:
            case "circle":
                return "A round shape"
            case "square":
                return "A shape with four equal sides"
            case "triangle":
                return "A shape with three sides"
            case _:
                return "Unknown shape"
\end{lstlisting}

Detected refactoring:
\begin{itemize}[leftmargin=*]
  \item \code{REPLACE\_CONDITIONAL\_WITH\_PATTERN\_MATCHING} in function \code{describe} of class \code{ShapeProcessor}
\end{itemize}

\subsection{Replace Conditional With Pattern Matching --- Guard Clauses}

When at least one \code{case} branch uses an \code{if} guard (e.g., \code{case x if x < 0:}), the resulting code mixes literal patterns and guard-based dispatch. RefactoringMiner detects the refactoring as a single event and does not report the guard conditions as additional refactorings.

\begin{lstlisting}[caption={Before: range checks in an \texttt{if}/\texttt{elif}/\texttt{else} chain.}]
class Classifier:
    def classify(self, value):
        if value < 0:
            return "negative"
        elif value == 0:
            return "zero"
        elif value < 100:
            return "small positive"
        else:
            return "large positive"
\end{lstlisting}

\begin{lstlisting}[caption={After: \texttt{match} with guard clauses.}]
class Classifier:
    def classify(self, value):
        match value:
            case x if x < 0:
                return "negative"
            case 0:
                return "zero"
            case x if x < 100:
                return "small positive"
            case _:
                return "large positive"
\end{lstlisting}

Detected refactoring:
\begin{itemize}[leftmargin=*]
  \item \code{REPLACE\_CONDITIONAL\_WITH\_PATTERN\_MATCHING} in function \code{classify} of class \code{Classifier}
\end{itemize}

\subsection{Replace Pattern Matching With Conditional --- Inverse Direction}

The same transformation in reverse. Useful for projects that revert to conditionals for compatibility with Python~3.9 or for readability preferences.

\begin{lstlisting}[caption={Before: \texttt{match}/\texttt{case} dispatch.}]
class ShapeProcessor:
    def describe(self, shape):
        match shape:
            case "circle":
                return "A round shape"
            case "square":
                return "A shape with four equal sides"
            case "triangle":
                return "A shape with three sides"
            case _:
                return "Unknown shape"
\end{lstlisting}

\begin{lstlisting}[caption={After: reverted to \texttt{if}/\texttt{elif}/\texttt{else}.}]
class ShapeProcessor:
    def describe(self, shape):
        if shape == "circle":
            return "A round shape"
        elif shape == "square":
            return "A shape with four equal sides"
        elif shape == "triangle":
            return "A shape with three sides"
        else:
            return "Unknown shape"
\end{lstlisting}

Detected refactoring:
\begin{itemize}[leftmargin=*]
  \item \code{REPLACE\_PATTERN\_MATCHING\_WITH\_CONDITIONAL} in function \code{describe} of class \code{ShapeProcessor}
\end{itemize}

% -----------------------------------------------------------------------
\section{Design and Implementation Notes}

\subsection{Type Annotation Detection}

Detection happens at three distinct AST levels, each handled by a dedicated diff class:

\begin{enumerate}[leftmargin=*, label=\arabic*.]
  \item \textbf{\code{UMLOperationDiff} --- parameters and return type.}
        The diff compares matched \code{UMLParameter} objects from the before/after versions of a function and inspects the \code{hasTypeAnnotation} and \code{hasExplicitReturnTypeAnnotation} flags, which are propagated by the Python AST builder (\code{PyDeclarationASTBuilder}).

  \item \textbf{\code{UMLAttributeDiff} --- class-level attributes.}
        The diff inspects the annotation strings of \code{UMLAttribute} objects from both versions.

  \item \textbf{\code{UMLOperationBodyMapper} --- local variables.}
        Variable declarations inside function bodies are compared via \code{VariableDeclaration} objects in the \code{VariableReplacementAnalysis} pass.
\end{enumerate}

Four new refactoring record classes carry the results:

\begin{itemize}[leftmargin=*]
  \item \textbf{\code{AddTypeAnnotationRefactoring}} --- parameterised by \code{RefactoringType} to cover both add and remove cases for parameters and return types.
  \item \textbf{\code{AddVariableTypeAnnotationRefactoring}} --- covers local variables and class attributes.
  \item \textbf{\code{ChangeTypeAnnotationRefactoring}} --- captures \code{typeBefore} and \code{typeAfter} for function-level changes.
  \item \textbf{\code{ChangeVariableTypeAnnotationRefactoring}} --- captures \code{typeBefore} and \code{typeAfter} for variable-level changes.
\end{itemize}

An important ordering constraint: the CHANGE variants are checked \emph{before} the ADD/REMOVE variants. This prevents a \code{List[int]} $\to$ \code{list[int]} change from being double-reported as a remove followed by an add.

\subsection{Structural Pattern Matching Detection}

The Python AST builder already represents \code{match}/\code{case} as \code{Code\-Element\-Type.\-SWITCH\_STATEMENT} (via \code{LangSwitchStatement}), matching its conceptual equivalence with Java's \code{switch}. The \code{elif} branches are represented as nested \code{LangIfStatement} nodes, giving \code{CodeElementType.IF\_STATEMENT}.

Detection runs inside \code{UMLOperationBodyMapper} as two new private methods:

\begin{itemize}[leftmargin=*]
  \item \textbf{\code{detectIfToPatternMatch}.} Scans unmatched composite nodes from version~1 (\code{IF\_STATEMENT}) and version~2 (\code{SWITCH\_STATEMENT}). For each candidate pair, it applies a \emph{subject-in-condition heuristic}: it checks whether the subject expression of the \code{match} statement (the first expression of the \code{SWITCH} node) appears as a substring of any condition expression of the \code{if} statement. Nested \code{elif} branches are excluded by skipping any \code{IF\_STATEMENT} whose parent is also an \code{IF\_STATEMENT}.

  \item \textbf{\code{detectPatternMatchToIf}.} The exact mirror: scans \code{SWITCH} in version~1 and outermost \code{IF} in version~2, applying the same heuristic.
\end{itemize}

Both methods are called from two sites in \code{UMLOperationBodyMapper} to cover the two code paths that Python method comparisons can traverse: the direct \code{processInnerNodes} path used by the \code{UMLOperation} constructor, and the \code{processCompositeStatements} path.

% -----------------------------------------------------------------------
\section{Files Changed}

\subsection{New Production Files}

\begin{table}[ht]
  \centering
  \caption{New source files introduced by these changes.}
  \small
  \begin{tabular}{@{}ll@{}}
    \toprule
    \textbf{Class} & \textbf{Purpose} \\
    \midrule
    \code{AddTypeAnnotationRefactoring}               & Add/remove parameter and return annotations \\
    \code{AddVariableTypeAnnotationRefactoring}        & Add/remove variable annotations \\
    \code{ChangeTypeAnnotationRefactoring}             & Changed parameter/return annotations \\
    \code{ChangeVariableTypeAnnotationRefactoring}     & Changed variable annotations \\
    \code{ReplaceConditionalWithPatternMatchingRefactoring} & Refactoring record for \code{if} $\to$ \code{match} \\
    \code{ReplacePatternMatchingWithConditionalRefactoring} & Refactoring record for \code{match} $\to$ \code{if} \\
    \bottomrule
  \end{tabular}
\end{table}

\subsection{Modified Production Files}

\begin{table}[ht]
  \centering
  \caption{Existing files modified by these changes.}
  \small
  \begin{tabular}{@{}ll@{}}
    \toprule
    \textbf{File} & \textbf{Change} \\
    \midrule
    \code{RefactoringType}              & 11 new enum constants and entries in the \code{ALL} array \\
    \code{UMLOperationBodyMapper}       & Pattern matching detection; local-variable annotation detection \\
    \code{UMLOperationDiff}             & Parameter and return annotation add/change/remove detection \\
    \code{UMLAttributeDiff}             & Class attribute annotation add/change/remove detection \\
    \code{UMLOperation}                 & \code{hasExplicitReturnTypeAnnotation} flag \\
    \code{UMLParameter}                 & \code{hasTypeAnnotation} flag \\
    \code{UMLModelAdapter}              & Propagates annotation flags from AST nodes \\
    \code{LangMethodDeclaration}        & Stores return annotation presence \\
    \code{LangSingleVariableDeclaration}& Stores parameter annotation presence \\
    \code{PyDeclarationASTBuilder}      & Reads annotation presence from the Python parse tree \\
    \bottomrule
  \end{tabular}
\end{table}

\subsection{New Test Files}

Eight new JUnit~5 test classes were added, covering approximately 30 scenarios:

\begin{itemize}[leftmargin=*, itemsep=1pt]
  \item \code{TestAddTypeAnnotationRefactoring}
  \item \code{TestAddVariableTypeAnnotationRefactoring}
  \item \code{TestRemoveFunctionTypeAnnotationRefactoring}
  \item \code{TestRemoveVariableTypeAnnotationRefactoring}
  \item \code{TestChangeFunctionTypeAnnotationRefactoring}
  \item \code{TestChangeVariableTypeAnnotationRefactoring}
  \item \code{TestReplaceConditionalWithPatternMatchingRefactoring}
  \item \code{TestReplacePatternMatchingWithConditionalRefactoring}
\end{itemize}

\end{document}
